\documentclass[10pt,a4paper]{amsart}
\usepackage{amsmath,amssymb,amsthm,graphicx}
\usepackage{verbatim}
\newenvironment{reference}{\comment}{\endcomment}
%\newenvironment{reference}{}{}
\newenvironment{slogan}{\comment}{\endcomment}
\newenvironment{history}{\comment}{\endcomment}

% For commutative diagrams we use Xy-pic
\usepackage[all]{xy}

% We use 2cell for 2-commutative diagrams.
\xyoption{2cell}
% Theorem environments.
%
\theoremstyle{plain}
\newtheorem{theorem}[subsection]{Theorem}
\newtheorem{proposition}[subsection]{Proposition}
\newtheorem{lemma}[subsection]{Lemma}

\theoremstyle{definition}
\newtheorem{definition}[subsection]{Definition}
\newtheorem{example}[subsection]{Example}
\newtheorem{exercise}[subsection]{Exercise}
\newtheorem{situation}[subsection]{Situation}

\theoremstyle{remark}
\newtheorem{remark}[subsection]{Remark}
\newtheorem{remarks}[subsection]{Remarks}

\numberwithin{equation}{subsection}

% Macros
%
\def\lim{\mathop{\mathrm{lim}}\nolimits}
\def\colim{\mathop{\mathrm{colim}}\nolimits}
\def\Spec{\mathop{\mathrm{Spec}}}
\def\Hom{\mathop{\mathrm{Hom}}\nolimits}
\def\Ext{\mathop{\mathrm{Ext}}\nolimits}
\def\SheafHom{\mathop{\mathcal{H}\!\mathit{om}}\nolimits}
\def\SheafExt{\mathop{\mathcal{E}\!\mathit{xt}}\nolimits}
\def\Sch{\mathit{Sch}}
\def\Mor{\mathop{\mathrm{Mor}}\nolimits}
\def\Ob{\mathop{\mathrm{Ob}}\nolimits}
\def\Sh{\mathop{\mathit{Sh}}\nolimits}
\def\NL{\mathop{N\!L}\nolimits}
\def\CH{\mathop{\mathrm{CH}}\nolimits}
\def\proetale{{pro\text{-}\acute{e}tale}}
\def\etale{{\acute{e}tale}}
\def\QCoh{\mathit{QCoh}}
\def\Ker{\mathop{\mathrm{Ker}}}
\def\Im{\mathop{\mathrm{Im}}}
\def\Coker{\mathop{\mathrm{Coker}}}
\def\Coim{\mathop{\mathrm{Coim}}}

% Boxtimes
%
\DeclareMathSymbol{\boxtimes}{\mathbin}{AMSa}{"02}

%
% Macros for moduli stacks/spaces
%
\def\QCohstack{\mathcal{QC}\!\mathit{oh}}
\def\Cohstack{\mathcal{C}\!\mathit{oh}}
\def\Spacesstack{\mathcal{S}\!\mathit{paces}}
\def\Quotfunctor{\mathrm{Quot}}
\def\Hilbfunctor{\mathrm{Hilb}}
\def\Curvesstack{\mathcal{C}\!\mathit{urves}}
\def\Polarizedstack{\mathcal{P}\!\mathit{olarized}}
\def\Complexesstack{\mathcal{C}\!\mathit{omplexes}}
% \Pic is the operator that assigns to X its picard group, usage \Pic(X)
% \Picardstack_{X/B} denotes the Picard stack of X over B
% \Picardfunctor_{X/B} denotes the Picard functor of X over B
\def\Pic{\mathop{\mathrm{Pic}}\nolimits}
\def\Picardstack{\mathcal{P}\!\mathit{ic}}
\def\Picardfunctor{\mathrm{Pic}}
\def\Deformationcategory{\mathcal{D}\!\mathit{ef}}
\usepackage{cleveref}
\crefname{originalreference}{original reference}{original references}
\Crefname{originalreference}{Original reference}{Original references}
\makeatletter
\newlabel{chow-lemma-rational-function}{{02RQ}{source}{}{sourcepointer-1}{}}
\newlabel{chow-lemma-rational-function@cref}{{[originalreference][0][]02RQ}{[1][0][]source}{}{}{}}
\makeatother
\crefname{originalreference}{original reference}{original references}
\Crefname{originalreference}{Original reference}{Original references}
\makeatletter
\newlabel{chow-lemma-proper-pushforward-alteration}{{02RT}{source}{}{sourcepointer-2}{}}
\newlabel{chow-lemma-proper-pushforward-alteration@cref}{{[originalreference][0][]02RT}{[1][0][]source}{}{}{}}
\makeatother
\title{Principal divisors generate rational equivalence on a variety}
\author{The Stacks Project Authors}
\date{stacks-project:intersection.tex:local-snapshot-20261012; GNU FDL 1.2 or later}
\begin{document}
\maketitle
\noindent\textit{Editorial note.} All delivered mathematical text is the author's own, copied verbatim from the pinned Stacks Project chapter file intersection.tex; no statement or proof step was rewritten, reordered or removed.\par\medskip
\noindent\textit{Editorial note.} Span recomputation by label pairing (LANE\_BRIEF 10.2): in this chapter the author's proof environment follows the statement environment as a standalone block; the spec covers the complete environment pair (statement span, proof span) as recorded in the pool candidate.\par\medskip
\noindent\textit{Editorial note.} Occupation scan: combined claim\_key + statement\_anchor + generic-line scan over all registration files (active and withdrawn) found 0 hits for this chapter (is-occupied.json, 2026-10-14); consistent with the duty manager's six-check dispatch.\par\medskip
\ifdefined\hypertarget\hypertarget{sourcepointer-1}{}\fi
\noindent\textit{Source reference.} Original reference 02RQ: sources/stacks-chow/tags.txt, source lines 3633--3633. Cross-chapter standard prerequisite cited by the author inside the counted span; it is used as a pointer and its own proof is not part of this item.\par\medskip
\ifdefined\hypertarget\hypertarget{sourcepointer-2}{}\fi
\noindent\textit{Source reference.} Original reference 02RT: sources/stacks-chow/tags.txt, source lines 3636--3636. Cross-chapter standard prerequisite cited by the author inside the counted span; it is used as a pointer and its own proof is not part of this item.\par\medskip
\begin{lemma}
\label{lemma-rational-equivalence}
Let $X$ be a variety. Let $W \subset X$ be a subvariety
of dimension $k + 1$. Let $f \in \mathbf{C}(W)^*$ be a nonzero rational
function on $W$. Then $\text{div}_W(f)$ is rationally equivalent to zero on
$X$. Conversely, these principal divisors generate the abelian group of
cycles rationally equivalent to zero on $X$.
\end{lemma}

\begin{proof}
The first assertion follows from
Chow Homology, Lemma \ref{chow-lemma-rational-function}.
More precisely, let $W' \subset X \times \mathbf{P}^1$ be the closure
of the graph of $f$. Then $\text{div}_W(f) = [W'_0]_k - [W'_\infty]$
in $Z_k(W) \subset Z_k(X)$, see part (6) of
Chow Homology, Lemma \ref{chow-lemma-rational-function}.

\medskip\noindent
For the second, let $W' \subset X \times \mathbf{P}^1$ be a closed
subvariety of dimension $k + 1$ which dominates $\mathbf{P}^1$.
We will show that $[W'_0]_k - [W'_\infty]_k$ is a principal divisor
which will finish the proof. Let $W \subset X$ be the image of $W'$
under the projection to $X$. Then $W \subset X$ is a closed subvariety
and $W' \to W$ is proper and dominant with fibres of dimension $0$
or $1$. If $\dim(W) = k$, then $W' = W \times \mathbf{P}^1$ and we
see that $[W'_0]_k - [W'_\infty]_k = [W] - [W] = 0$.
If $\dim(W) = k + 1$, then
$W' \to W$ is generically finite\footnote{If $W' \to W$ is birational,
then the result follows from
Chow Homology, Lemma \ref{chow-lemma-rational-function}.
Our task is to show that even if $W' \to W$
has degree $> 1$ the basic rational equivalence
$[W'_0]_k \sim_{rat} [W'_\infty]_k$ comes from a principal divisor
on a subvariety of $X$.}. Let $f$ denote the projection $W' \to \mathbf{P}^1$
viewed as an element of $\mathbf{C}(W')^*$. Let
$g = \text{Nm}(f) \in \mathbf{C}(W)^*$ be the norm. By
Chow Homology, Lemma \ref{chow-lemma-proper-pushforward-alteration}
we have
$$
\text{div}_W(g) = \text{pr}_{X, *}\text{div}_{W'}(f)
$$
Since $\text{div}_{W'}(f) = [W'_0]_k - [W'_\infty]_k$ 
by Chow Homology, Lemma \ref{chow-lemma-rational-function}
the proof is complete.
\end{proof}

\end{document}
